\documentclass[fontsize=10pt]{scrartcl}

\usepackage[columns=3]{cheatsheet} % layout, colors and macros, see cheatsheet.sty
\usepackage[ngerman]{babel}        % language of the content (hyphenation, date)

\title{Template}
\author{Nicolas Huber}
\date{\today}

%%%%%%%%%%%%%%%%%%%%%%%%%%%%%%%%%%%%%%%%%%%%%%%%%%%%%%%%%%
%                     KEEP IN MIND
% "|" use \mid for relations (not a delimiter)
% "|" use \vert in the context of an ordinal (not a delimiter)
% "|" use \lvert and \rvert (or \abs) as delimiters (not operators)

\begin{document}
\begin{cheatsheet}
\maketitle

\section{Titel}
\subsection{Untertitel}
\subsubsection{Unter-Untertitel}

\subsubsection{Buchstaben}
\begin{itemize}
    \item Griechische Buchstaben (klein): \\
          $\alpha \beta \gamma \delta \epsilon \varepsilon \zeta \eta \theta \vartheta \iota \kappa \lambda \mu \nu \xi
           \pi \varpi \rho \varrho \sigma \varsigma \tau \upsilon \phi \varphi \chi \psi \omega$
    \item Griechische Buchstaben (gross): \\
          $A B \Gamma \Delta E Z H \Theta I K \Lambda M N \Xi O \Pi P \Sigma T \Upsilon \Phi X \Psi \Omega$
    \item Kalligraphische Symbole: \\ $\mathcal{ABCDEFGHIJKLMNOPQRSTUVWXYZ}$
\end{itemize}

\subsubsection{Notation}
\begin{itemize}
    \item Definitionen: \defn Satz des Pythagoras
    \item Bemerkungen: \bem Tipps zum Faktorisieren
    \item Formeln: \eqbox{x_{1,2} = \frac{-b \pm \sqrt{b^2-4ac}}{2a}}
    \item Beweise: $\square$, $\blacksquare$
    \item Index: $x_1$
    \item Bruch (abgesetzt): \[ f(x) = \frac{1}{2} x^2 = \dfrac{1}{2} x^2 = \tfrac{1}{2} x^2 \]
    \item Bruch (im Text): $f(x) = \frac{1}{2} x^2 = \dfrac{1}{2} x^2 = \tfrac{1}{2} x^2$
\end{itemize}

\subsubsection{Pfeile}
$\leftarrow, \Leftarrow, \rightarrow, \Rightarrow, \leftrightarrow, \Longleftrightarrow, \mapsto, \longmapsto$

\subsubsection{Zahlenbereiche}
\begin{symbols}
    \mathbb{N}     & natürliche Zahlen \\
    \mathbb{N}_{0} & natürliche Zahlen mit 0 \\
    \mathbb{Z}     & ganze Zahlen \\
    \mathbb{Q}     & rationale Zahlen \\
    \mathbb{R}     & reelle Zahlen \\
    \mathbb{C}     & komplexe Zahlen \\
\end{symbols}

\subsubsection{Gleichheit}
\begin{symbols}
    =             & gleich \\
    \neq          & ungleich \\
    \coloneqq     & nach Definition gleich, Wertzuweisung \\
    \approx       & ungefähr gleich \\
    \equiv        & identisch \\
    <, \leq       & kleiner bzw. kleiner gleich \\
    >, \geq       & grösser bzw. grösser gleich \\
    \ll           & viel kleiner \\
    \gg           & viel grösser \\
\end{symbols}

\subsubsection{Mengen}
\begin{symbols}
    \emptyset, \{\}         & leere Menge \\
    \in, \ni, \notin        & Element von, enthält, nicht Element von \\
    \cup, \bigcup           & vereinigt mit, Vereinigung \\
    \cap, \bigcap           & geschnitten mit, Schnittmenge \\
    \setminus               & Differenz (ohne) \\
    A^{c}, \overline{A}     & Komplement \\
    \subset, \supset        & echte Teilmenge bzw. Obermenge \\
    \subseteq, \supseteq    & Teilmenge bzw. Obermenge \\
\end{symbols}

\subsubsection{Intervalle}
\begin{itemize}
    \item offen: $(a,b) = \mathopen]a,b\mathclose[ = \{x \in \mathbb{R} \mid a < x < b\}$
    \item abgeschlossen: $[a,b] = \{x \in \mathbb{R} \mid a \leq x \leq b\}$
\end{itemize}

\subsubsection{Logik}
\begin{symbols}
    \land            & logisches UND \\
    \lor             & logisches ODER \\
    \neg             & Negation \\
    \Rightarrow      & Implikation (daraus folgt) \\
    \iff             & Äquivalenz (genau dann, wenn) \\
    \because         & weil \\
    \therefore       & deshalb \\
    \forall          & Allquantor (für alle) \\
    \exists          & Existenzquantor (es existiert) \\
\end{symbols}

\subsubsection{Funktionen}
\begin{itemize}
    \item Trigonometrische Funktionen: $\sin x, \cos x, \tan x$
    \item Wurzeln: $\sqrt[n]{x}$, Quadratwurzel: $\sqrt{x}$
    \item Exponentialfunktion: $f(x) = e^{2x}$
    \item Logarithmusfunktion: $f(x) = \log_2(x)$
    \item Floor- und Ceiling-Funktion: $\floor{x}, \ceil{x}$
    \item Betragsfunktion: $f(x) = \abs{x}$
\end{itemize}

\subsubsection{Summen, Produkte und Grenzwerte}
\begin{itemize}
    \item Summe: $\sum_{i=1}^n i = \frac{n(n+1)}{2}$
    \item Summe (mit \texttt{\textbackslash limits}): $\sum\limits_{i=1}^n i = \frac{n(n+1)}{2}$
    \item Produkt: $\prod_{i=1}^n i = n!$
    \item Produkt (mit \texttt{\textbackslash limits}): $\prod\limits_{i=1}^n i = n!$
    \item Maximum: $\max\limits_{x \in (a,b)} f(x)$, Minimum: $\min\limits_{x \in (a,b)} f(x)$
    \item Limes: $\lim\limits_{n \to \infty} \frac{1}{n} = 0$
\end{itemize}

\subsubsection{Integral- und Differentialrechnung}
\begin{itemize}
    \item Ableitung: $f(x) = 3x^3,\ f'(x) = 9x^2$
    \item Integral: $\int x^2 \dif x = \frac{x^3}{3} + C,\ C \in \mathbb{R}$
    \item Verknüpfung: $(f \circ g)(x) = f(g(x))$
    \item Differentialoperatoren: $\nabla, \partial$
\end{itemize}

\subsubsection{Vektoren und Matrizen}
\begin{itemize}
    \item Vektor: $\vec{v} = \vect{v} = \matp{1 \\ 2 \\ 3} = \matp{1 & 2 & 3}^T$
    \item Matrix mit runden Klammern: $A = \matp{1 & 2 & 3 \\ a & b & c}$
    \item Matrix mit eckigen Klammern: $A = \matb{1 & 2 & 3 \\ a & b & c}$
    \item Matrix ohne Klammern: $A = \mat{1 & 2 & 3 \\ a & b & c}$
    \item Determinante: $\det(A) = \begin{vmatrix} 2 & 1 \\ -2 & 1 \end{vmatrix} = 4$
    \item Norm und Betrag: $\norm{a \vec{v}} = \abs{a} \, \norm{\vec{v}}$
\end{itemize}

\subsubsection{Diverses}
\begin{symbols}
    \infty           & unendlich \\
    \propto          & proportional zu \\
    a \bmod n        & Modulo \\
    \times           & mal (Kreuzprodukt) \\
    \cdot            & mal (Skalarprodukt) \\
    a \mid b         & $a$ teilt $b$ \\
    \Re, \Im         & Real- und Imaginärteil \\
    \binom{n}{k}     & Binomialkoeffizient \\
\end{symbols}

\end{cheatsheet}
\end{document}
